\documentclass[10pt,a4paper]{amsart}
\usepackage[margin=19mm]{geometry}
\usepackage{amsmath,amssymb,mathtools}
\usepackage[scr=rsfs,cal=euler]{mathalfa}
\usepackage{xfrac}
\usepackage[unicode,hidelinks,bookmarks=false]{hyperref}
\newcommand{\ie}{i.e.\ }
\newcommand{\cf}{cf.\ }
\newcommand{\resp}{respectively\ }
\newcommand{\Subsection}{\subsection}
\providecommand{\nobreakdash}{\nobreak\hbox{-}}
\setlength{\emergencystretch}{2em}
\multlinegap0pt
\usepackage{mathtools}
\usepackage{enumitem}
\usepackage{tikz}
\newtheorem{theorem}{Theorem}[section]
\newtheorem{proposition}[theorem]{Proposition}
\newtheorem{lemma}[theorem]{Lemma}
\newtheorem{corollary}[theorem]{Corollary}
\newtheorem{definition}[theorem]{Definition}
\newtheorem{example}[theorem]{Example}
\newtheorem{remark}[theorem]{Remark}
\newcommand{\VR}{\mathrm{VR}}
\newcommand{\Cech}{\check{C}}
\newcommand{\Dow}{\mathrm{D}}
\newcommand{\DR}{\mathrm{DR}}
\newcommand{\DRR}{\mathrm{DRR}}
\newcommand{\Flag}{\mathrm{F}}
\newcommand{\St}{\mathrm{St}}
\newcommand{\Lk}{\mathrm{Lk}}
\newcommand{\PH}{\mathrm{PH}}
\newcommand{\R}{\mathbb{R}}
\begin{document}
\title[Geometry-Induced Hodge Stars on Rips and Dowker--Rips Comple]{Geometry-Induced Hodge Stars on Rips and Dowker--Rips Complexes}
\author{Jiahui Chen; Sebastian Wilcox}
\date{}
\maketitle
\noindent\emph{Source and licence.} Geometry-Induced Hodge Stars on Rips and Dowker--Rips Complexes. Version arXiv:2607.18692v1 (CC BY 4.0). This material is distributed under the Creative Commons Attribution 4.0 International licence, http://creativecommons.org/licenses/by/4.0/, as recorded in the arXiv OAI licence metadata for this exact version and shown on the abstract page https://arxiv.org/abs/2607.18692v1. Full rights notice: licenses/arxiv-2607.18692-CC-BY-4.0.txt.\par\smallskip
\noindent\textit{Editorial scope.} Counted unit: the original \texttt{proposition} environment \texttt{prop:witness-separation} at source lines 396--414 together with its complete proof at lines 416--432; the delivered document keeps source lines 381--433 so that every referenced label and every citation resolves inside it. Native editable source is the paper's own concatenated TeX; the AMS wrapper is editorial formatting only. No publisher class, upstream script or unrelated artwork is redistributed.
\section{Theoretical Properties}
\label{sec:theory}

Throughout, self-adjointness is taken with respect to the weighted inner product. Thus an operator
$T:C^k(K)\to C^k(K)$ is self-adjoint if
\[
\langle T\alpha,\beta\rangle_{k,W}
=
\langle\alpha,T\beta\rangle_{k,W}
\qquad
\text{for all }\alpha,\beta\in C^k(K),
\]
equivalently if $W_kT$ is symmetric. Note that $\Delta_k^W$ is generally not symmetric in the
standard basis, so this qualifier is essential.


\begin{proposition}[Asymptotic soft-witness decay rate]
\label{prop:witness-separation}
Let $X,Y$ be finite metric spaces in a common ambient metric space. Define
$s_t(\sigma;Y)=\sum_{y\in Y}\exp(-t^{-1}\sum_{x\in\sigma}d(x,y)^2)$, set
$D(\sigma,y)=\sum_{x\in\sigma}d(x,y)^2$, and let
$D^{*}(\sigma)=\min_{y\in Y}D(\sigma,y)$.
Then
\[
\lim_{t\to0^{+}}\bigl(-t\log s_t(\sigma;Y)\bigr)=D^{*}(\sigma).
\]
Moreover, for the metric relation $R_\epsilon=\{(x,y):d(x,y)\le\epsilon\}$:
\begin{enumerate}[leftmargin=2em,label=(\roman*)]
\item if $\sigma\in\Dow_\epsilon(X,Y)$, then $D^{*}(\sigma)\le \#\sigma\,\epsilon^2$;
\item if $\sigma\in\DR_\epsilon(X,Y)\setminus\Dow_\epsilon(X,Y)$, then
$D^{*}(\sigma)>\epsilon^2$.
\end{enumerate}
Thus the soft witness support has a precise small-bandwidth decay rate, and flagified-only
simplices have decay rate bounded below by $\epsilon^2$.
\end{proposition}

\begin{proof}
Since $Y$ is finite, the largest exponential term controls the sum:
\[
e^{-D^{*}(\sigma)/t}\le s_t(\sigma;Y)\le |Y|e^{-D^{*}(\sigma)/t}.
\]
Taking $-t\log(\cdot)$ gives
\[
D^{*}(\sigma)-t\log |Y|\le -t\log s_t(\sigma;Y)\le D^{*}(\sigma),
\]
and the limit follows. If $\sigma\in\Dow_\epsilon(X,Y)$, there exists a common witness $y_0$
with $d(x,y_0)\le\epsilon$ for every $x\in\sigma$, so
$D^{*}(\sigma)\le D(\sigma,y_0)\le \#\sigma\,\epsilon^2$. If
$\sigma\in\DR_\epsilon(X,Y)\setminus\Dow_\epsilon(X,Y)$, then no $y\in Y$ witnesses all vertices.
For each $y$, some $x\in\sigma$ satisfies $d(x,y)>\epsilon$, hence
$D(\sigma,y)>\epsilon^2$. Taking the minimum over finite $Y$ gives
$D^{*}(\sigma)>\epsilon^2$.
\end{proof}


\end{document}
